% Part: lambda-calculus
% Chapter: introduction
% Section: lambda-definability

\documentclass[../../../include/open-logic-section]{subfiles}

\begin{document}

\olfileid{lam}{int}{rep}
\olsection{\usetoken{S}{lambda definable} अंकगणितीय फलने}

लॅम्डा कलन संगणनाचे प्रतिमान कसे ठरू शकते? सुरुवातीला तर या
विधानाचा अर्थ काय, हेही स्पष्ट नाही. नैसर्गिक संख्यांवरील
संगणनक्षमतेची चर्चा करायची असेल, तर त्या संख्यांचे योग्य
निरूपण शोधावे लागेल. आश्चर्यकारकरीत्या चांगले काम करणारे
एक निरूपण पुढे दिले आहे.

\begin{defn}
प्रत्येक नैसर्गिक संख्या~$n$ साठी \emph{चर्च अंक}
$\num{n}$ हे लॅम्डा पद $\lambd[x][\lambd[y][(x(x(x(\dots
x(y)))))]]$ असे परिभाषित करा; यात एकूण $n$ वेळा $x$ येते.
\end{defn}

$\num{n}$ ही पदे “पुनरावर्तक” आहेत: $f$ हा निविष्ट दिल्यास
$\num{n}$ हे $y$ ला $f^n(y)$ शी जोडणारे फलन परत करते.
प्रत्येक अंक सामान्य रूपात आहे, हे लक्षात घ्या. आता
एखादे लॅम्डा पद नैसर्गिक संख्यांवरील फलनाचे “संगणन करते”
याचा अर्थ सांगता येईल.

\begin{defn}
$f(x_0, \dots, x_{k-1})$ हे $\Nat$ पासून $\Nat$ मध्ये
जाणारे $k$-स्थानी आंशिक फलन असो. एखादे $\lambd$-पद~$F$
हे~$f$ ला \emph{!!{lambda define}s} असे तेव्हाच म्हणतो, जेव्हा
नैसर्गिक संख्यांच्या प्रत्येक $n_0$, \dots,~$n_{k-1}$
अनुक्रमासाठी, $f(n_0, \dots, n_{k-1})$ परिभाषित असेल तर
\[
F\, \num{n_0}\, \num{n_1} \dots \num{n_{k-1}} \red \num{f(n_0, n_1, \dots,
  n_{k-1})}
\]
असे असते; आणि तसे नसेल तर $F\, \num{n_0}\, \num{n_1}
\dots \num{n_{k-1}}$ याला सामान्य रूप नसते.
\end{defn}

\begin{thm}
\ollabel{thm:lambda-def}
फलन $f$ हे आंशिक संगणनक्षम फलन असेल तेव्हाच आणि तेव्हाच
एखाद्या लॅम्डा पदाने ते !!{lambda defined} असते.
\end{thm}

\begin{explain}
हे प्रमेय काहीसे थक्क करणारे आहे. संगणनाचे प्रतिमान म्हणून
लॅम्डा कलन अगदी साधे कलन आहे; त्यात लॅम्डा अमूर्तीकरण आणि
उपयोजन एवढ्याच क्रिया आहेत!{} तरी या मोजक्या साधनांतून
कोणतीही संगणनप्रक्रिया साकार करता येते.
\end{explain}

\end{document}
